\chapter{Aplicación de control básico}
Este capítulo presenta el primer hito funcional del proyecto, una aplicación Flutter Web capaz de conectarse al dron Ardupilot, armar sus motores, despegar, navegar en las seis direcciones básicas y aterrizar. Esta versión establece la arquitectura base y el modelo de gestión de estados sobre los que se basarán todas las arquitecturas posteriores. El apartado 5.1 describe la arquitectura general del sistema. El 5.2 detalla la cadena de comandos seguida desde la interfaz del usuario hasta el dron. El 5.3 explica la máquina de estado del dron y su gestión mediante un Provider. El 5.4 presenta las pantallas de la interfaz. El último apartado, 5.5, describe el proceso de validación mediante SITL.

\section{Arquitectura del sistema}
La primera versión está compuesta por tres componentes que se comunican entre sí mediante dos canales independientes. El frontend Flutter, el bakcend Dart y la estación tierra en Python.

El frontend Flutter Web es una aplicación que se ejecuta en el navegador y se estructura en torno a un gestor de estados centralizados, la clase \texttt{Provider}. Se comunica con el backend Dart mediante peticiones HTTP, peticiones \texttt{POST} para emitir comandos de vuelo y peticiones \texttt{GET} periódicas para obtener el estado actual del dron y los datos de telemetría (el polling mencionado anteriormente en la sección 3.5).
El backend Dart es un servidor HTTP construido con el framework Shelf, que escucha en el puerto 8080 y contiene una API REST. Recibe los comandos del frontend, los traduce en publicaciones MQTT hacia el broker HiveMQ, y mantiene actualizado al frontend sobre el estado del dron (frontend lo consulta mediante el polling).

La estación tierra Python es un script que se conecta al mismo broker que el servidor (broker.hivemq.com, puerto 1883), se suscribe al árbol de topics, interpreta los comandos recibidos y emite las instrucciones correspondientes al dron mediante la librería DronLink. Además de eso, publica confirmaciones de estado, y periódicamente la telemetría de vuela al servidor.
Esta arquitectura de tres componentes con el backend como intermediario es coherente con el modo de comunicación Global del DEE, descrito en el Capítulo 2. 

\section{Cadena de comandos}
Los comandos usados se mandan desde la interfaz Flutter hasta el dron a través de una pipeline de dos tramos. El primero, se trata de un tramo HTTP que comunica el frontend y el backend, el segundo, es un tramo MQTT que junta el backend con la estación tierra.

Cuando el usuario pulsa un botón en la interfaz, el DronProvider invoca el método correspondiente de la clase HttpLogic. Esta clase centraliza todas las llamadas HTTP y usa la librería http de Dart. Para los comandos sin parámetros (como son Connect, Arm, Land, Rtl o Disconnect), se utiliza el método privado sendPostRequest(url), que emite un POST y verifica que la respuesta sea 200 (OK). Para los comandos con parámetro (como takeoff o speed), se serializa el valor en la payload de la petición como JSON. Finalmente, para recibir, se usa el método getStatus( ), que realiza un GET al endpoint /status y devuelve una payload con los datos.

El backend, define un router Shelf con un handler. Al recibir una petición, cada handler actualiza el estado interno según corresponda, y publica el topic MQTT equivalente hacia la estación tierra. Para poner un ejemplo, cuando se envía el comando de armar, el handler de /arm, pone los flags de dronFlying y dronArmed en false como medida de precaución antes de publicar el /arm. O por ejemplo, el handler de /takeoff lee el campo altitude de la payload JSON y lo publica en /takeoff. 

Todos los handlers devuelven un código 200 (junto a las cabeceras CORS necesarias para permitir las peticiones desde el navegador) al frontend.

La estación tierra, suscrita a mobileFlutter/groundStation/\#, recibe todos los mensajes que publica el backend, y llama a los métodos de Dronlink correspondientes (como podría ser /arm - dron.arm( ), /takeoff - dron.tafeOff(altitude) o /land - dron.Land( ) ). Una vez la operación se ha realizado, la estación tierra publica la confirmación en el topic de respuesta (groundStation/mobileFlutter/armed,flying,landed…).

El camino de retorno de la telemetría sigue el modelo de polling. La estación tierra llama periódicamente a dron.send\_telemetry\_info( ), que construye un diccionario con todos los datos de telemetría, lo serializa en formato JSON y lo publica en groundStation/mobileFlutter/telemetry, donde el backend los recibirá y almacenará. Finalmente, en el siguiente ciclo de polling, el frontend obtendrá dichos valores mediante un GET a /status.

\section{Máquina de estados}
El estado del dron se rastrea de forma independiente en los tres componentes de la aplicación. La fuente de mayor autoridad es la estación tierra, ya que tiene acceso directo al autopiloto del dron. El backend mantiene una réplica de este, actualizado mediantes los topics MQTT que envía la estación tierra. El frontend tiene su propia representación de los estados mediante el DronProvider.

Las variables de estado que gestiona el Provider son principalmente tres booleanos: isConnected, que pasa a verdadero cuando la estación tierra confirma la conexión con el dron, isArmed, que pasa a verdadero cuando se ha confirmado que los motores están armados y listos para el despegue, e isFlying, que permanece en verdadero desde la confirmación del despegue hasta la confirmación del aterrizaje. Estas variables son las que tienen el control de qué botones están disponibles en la interfaz. Por ejemplo, “Takeoff” solo está habilitado cuando isConnected e isArmed són verdaderos, o “Land” y “RTL” solo están habilitados cuando isFlying es true.

Un reto que ha sido destacable, es la confirmación del armado. ArduPilot ejecuta una secuencia de comprobaciones previas al armado, que pueden tardar varios segundos en completarse. El DronProvider se encarga de gestionar esta situación de la siguiente manera: Se implementa una deadline, cuando se pulsa ARM se activa un boolean “-waitingForArm” y empieza un contador de 5 segundos en “-armDeadline”. Si en uno de los ciclos de polling se detecta que armed = true, entonces se cancela la espera, pero si el plazo se termina sin la confirmación del armado, se entiende que no se ha podido armar y se muestra el mensaje “Not ready to arm”. Este mecanismo evita que la interfaz se desarme automáticamente o se quede en un estado de carga indefinido.

La telemetría recibida en cada ciclo de polling se almacena en una serie de variables de tipo double y String. Adicionalmente, cada posición GPS válida (distinta a 0 o null), se guarda en una lista “droneTrail”, para así representar la trayectoria del dron en el mapa.

\section{Interfaz de usuario}
La aplicación se compone de dos pantallas con navegación secuencial, implementadas como widgets sin estado (StatelessWidget) suscritos al DronProvider mediante “context.Watch<DronProvider>( ).

\subsection{Setup Screen}
La pantalla de configuración (SetupScreen) es el punto de partida de la aplicación. Contiene dos campos de texto para la configuración del vuelo, la altitud de despegue (con un rango válido entre 2,0 y 50,0 metros) y la velocidad de vuelo (con rango entre 1,0 y 15,0 metros por segundo). Ambos campos se validan en tiempo real gracias al Provider, que les atribuye una flag “isConfigValid” donde si se introduce un valor dentro de los márgenes se iguala a True. Si el valor está fuera del rango, el texto auxiliar se muestra en rojo, y el botón “START FLIGHT” permanece deshabilitado hasta que se introduzcan valores válidos.

La pantalla incluye tres botones de gestión: 
\begin{itemize}
  \item \textbf{Connect:} Llama a connectDron( ), que envía un POST a /connect, pone al provider en “isLoading” = True e inicia el polling.
  
\item \textbf{Start flight:} Navega a la pantalla de vuelo (mediante un Navigator.push, herramienta navegadora de Dart), y sólo está habilitado cuando isConnected e isConfigValid son True.

\item \textbf{Disconnect:}  Llama a diconnectDron( ), que envía un POST a /disconnect, detiene el polling y limpia todas las variables.
\end{itemize}

\subsection{Flight Screen}
La pantalla de vuelo (Flight Screen) es la interfaz operativa. Se organiza en zonas horizontales (rows).

La barra superior muestra la telemetría en tiempo real: altitud (ALT), velocidad (GS), rumbo magnético (HDG), latitud (LAT), longitud (LON) con 7 decimales (igual que en Mission Planner), estado del autopiloto (STATE), modo de vuelo (MODE) y nivel de batería. Este último, tiene un código de colores, donde el símbolo de la batería aparece verde por encima del 50%, naranja entre el 20% y 50% y rojo por debajo del 20%. En el centro de esta zona se encuentra un texto que muestra el “message” del Provider.

La zona central presenta un mapa satelital proporcionado por la librería flutter\_map con tiles de OpenStreetMap. El mapa muestra el icono del dron en su posición GPS actual, rotado según su rumbo “currentHeading”. Una flecha naranja superpuesta indica la dirección de vuelo, y la trazada del dron se indica con una polilínea púrpura “droneTrail”. La velocidad vectorial la indica una línea verde formada por el vector entre los campos de velocidad Vx y Vy, que aparte de mostrar la magnitud relativa de la velocidad, también indica la dirección del movimiento actual. Y finalmente una sombra circular girs que indica si el dron está volando.

Si el dispositivo del operador tiene señal GPS, se muestra su posición como un icono de móvil azul junto con un indicador de precisión codificado en colores (verde por debajo de 5 m, naranja por debajo de 20 m, rojo en caso superior).

A ambos lados del mapa se sitúan los controles de movimiento, implementados como un widget de ContinuousButton. Este widget utiliza la función propia de Dart GestureDetector para detectar onTapDown y onTapUp, de forma que mientras el botón esté pulsado, va llamando a DronProvider.startmove(direction), que envía un POST a /move con la payload JSON {direction: “Forward”} (por ejemplo). Al soltar el ContinuousButton, se llama a DronProvider.stopMove( ), que envía la dirección “Stop”, parando el movimiento del dron.
El widget tiene una animación al pulsarlo (usando AnimatedContainer), donde se reduce el tamaño y se cambia el color, para así indicar visualmente cuál se está pulsando. 
En esta primera versión, los controles de movimiento son flechas. En futuras versiones serán cambiadas por joysticks virtuales.

La barra inferior contiene cinco botones de vuelo principales: ARM, TAKEOFF, DISCONNECT, LAND y RTL. La habilitación de cada botón viene determinada por la combinación de “isLoading”, “isConnected”, “isArmed” e “isFlying”, garantizando que solo son accesibles los comandos válidos para cada estado del dron (es decir, que si está volando, “isFlying” = True, no se pueda desonnectar, o no pueda armarse si ya está armado, “isArmed” = True).

\section{Validación mediante SITL}

La aplicación base fue validada usando Mission Planner como emulador de vuelo en modo SITL (Software In The Loop). Esta modalidad permite ejecutar el firmware ArduPilot en el propio ordenador de desarrollo, simulando un vehículo con comportamiento realista sin necesidad de usar un dron real. La estación tierra Python se conecta al simulador mediante la cadena de conexión “tcp:127.0.0.1: 5763", que es el puerto TCP estándar que expone el SITL del ArduPilot.

La validación confirmó la correcta ejecución de la secuencia de vuelo completa: establecimiento de la conexión, armado de motores, despegue a la altitud configurada, navegación en las seis direcciones, retorno al punto de lanzamiento y aterrizaje. También permitió verificar el correcto funcionamiento del mecanismo de deadline del armado y la actualización continua de la telemetría en la interfaz durante todo el vuelo.
